\documentclass[11pt,a4paper]{article}
\usepackage{times}
\usepackage{latexsym}
\usepackage{xcolor}
\usepackage{url}
\usepackage{graphicx}
\usepackage{booktabs}
\usepackage{amsmath}
\usepackage[a4paper,left=1.25cm,right=1.25cm,top=2.0cm,bottom=2.0cm]{geometry}

\setlength{\emergencystretch}{3em}
\sloppy

% Underline-style emphasis: \emph{...} now renders as underlined text.
% normalem preserves italics for \textit{} and \em elsewhere in the doc.
\usepackage[normalem]{ulem}
\renewcommand{\emph}[1]{\uline{#1}}

\DeclareUrlCommand\code{%
  \urlstyle{tt}%
  \def\UrlBreaks{\do\_\do\.\do\/\do\(\do\)\do\-}%
  \def\UrlBigBreaks{\do\_\do\.\do\/}%
}

% Taxonomy / classification label macro --- italics for conceptual category names.
% Use \lbl{contradiction to evidence}, NOT \texttt{contradiction\_to\_evidence}.
\newcommand{\lbl}[1]{\textit{#1}}

\newcommand\BibTeX{B{\sc ib}\TeX}

\begin{document}

\title{A Comparative Study of Hallucinations in Open-Weight LLMs:\\
       NLI-Grounded Taxonomy Labeling Across Models, Datasets, and Prompt Variants}

\author{
  Hitansh Paliwal \\
  \texttt{paliwalhitansh@gmail.com}
  \and
  Team Member 2 (TODO) \\
  \texttt{TODO@email}
}

\date{}

\maketitle

\begin{abstract}
We present a comparative study of hallucination behavior in four open-weight instruction-tuned LLMs (Phi-4-mini, Mistral-7B, Qwen2.5-7B, Llama-3.1-8B) on HotpotQA and TruthfulQA under three prompt variants. Our contribution is a reproducible labeling pipeline that combines a DeBERTa-MNLI judge with score-based heuristics to assign each answer one of seven taxonomy labels (\lbl{supported}, \lbl{abstained}, \lbl{contradiction to evidence}, \lbl{entity error}, \lbl{attribute error}, \lbl{multi-hop reasoning error}, \lbl{unsupported inference}). A multi-signal supported-override rescues verbose-but-correct answers, and a cross-record pass flags \lbl{overconfident abstention failure}. Across 7{,}440 labeled records we find: (i) three of four models hallucinate at near-identical aggregate rates (87--92\%) on HotpotQA plain, while Llama-3.1 stands out at 40\%; (ii) prompt sensitivity is sharply non-uniform---the same reasoning prompt reduces Qwen/Mistral hallucinations by $\sim$17 points each but \emph{increases} Llama-3.1's by 17.8 points; (iii) under reasoning prompts, the dominant error type shifts from \lbl{contradiction to evidence} to \lbl{attribute error} (reaching 91.7\% for Qwen), suggesting reasoning trades contradictions for embedded numerical errors; (iv) on TruthfulQA, all four models converge near 83\% under reasoning, indicating misconception-style questions defeat prompt-based mitigation uniformly. We argue that LLM evaluation should report fine-grained type distributions and that prompt-level mitigation must be model-conditioned.
\end{abstract}

\section{Project Overview (10 points)}

\paragraph{} Large language models are now deployed for QA tasks where factual accuracy directly affects users. Despite their fluency, instruction-tuned LLMs continue to produce confident but factually incorrect outputs---\textit{hallucinations}. The problem is hard because hallucinations are coherent, plausible-but-wrong statements that surface fluency checks cannot detect. The model has no internal signal distinguishing what it knows from what it can fluently confabulate. We narrow this problem to: \textbf{given a question, a model answer, and human-annotated supporting facts, automatically classify whether the answer is hallucinated and what kind of hallucination it is}, producing one of seven taxonomy labels per record.

\paragraph{} Prior work falls into three families. \textit{Model-internal} methods (e.g., SelfCheckGPT) require white-box logits or multiple samples. \textit{Retrieval-grounded} pipelines (e.g., FactScore) decompose answers into atomic claims and verify each. \textit{NLI-based} judges compare a hypothesis (the answer) to a premise (gold passage). All produce a binary or graded verdict but few yield a fine-grained taxonomy of \textit{why} a hallucination occurred, and fewer apply a uniform pipeline across multiple LLMs. The gap we address is \textbf{a single, reproducible pipeline that classifies hallucinations into an interpretable taxonomy and applies uniformly across models, datasets, and prompts.}

\paragraph{} Our pipeline combines DeBERTa-MNLI with score-based heuristics. The NLI judge gives the supported/hallucinated decision and contributes the contradiction-probability score. For finer typing, three additional scores compete: a NER-based entity score, a number-mismatch score gated on entity-correctness, and an evidence-engagement score gated on question type. The argmax type wins; below a 0.5 confidence floor we return \lbl{unsupported inference}. Two safety nets handle NLI's known weaknesses: a multi-signal supported-override (token recall + embedding cosine + bidirectional NLI) for verbose-but-correct answers, and a cross-record second pass that flags overconfident-abstention failures by comparing plain and abstain prompt variants for the same question.

\paragraph{} We apply this pipeline uniformly to four open-weight LLMs (Phi-4-mini-instruct 3.8B, Mistral-7B-Instruct-v0.3, Qwen2.5-7B-Instruct, Llama-3.1-8B-Instruct) on 500 HotpotQA + 120 TruthfulQA questions $\times$ three prompts (\lbl{plain}, \lbl{abstain}, \lbl{reasoning}), totaling 7{,}440 labeled records. Primary metrics are EM, token-F1, hallucination rate, abstention rate, and the type distribution within hallucinated records.

\paragraph{} Three headline findings: (1) on HotpotQA \lbl{plain}, three of four models hallucinate at 87--92\%, while Llama-3.1 stands out at 40\%; (2) the same \lbl{reasoning} prompt reduces Mistral's and Qwen's hallucination rates by $\sim$17 points each but \emph{increases} Llama-3.1's by 17.8 points; (3) error \emph{types} shift under reasoning prompts: \lbl{contradiction to evidence} drops while \lbl{attribute error} rises (91.7\% for Qwen), and on TruthfulQA all four models converge near 83\% regardless of model.

\section{Ideas (10 points)}

\subsection{Idea 1: NLI-Based Evidence-Grounded Labeling}

HotpotQA includes human-annotated supporting facts---sentence-level pointers to gold passages---which makes the verification problem a textbook NLI task: does the premise (gold evidence) entail the hypothesis (model answer)? We use \texttt{microsoft/deberta-large-mnli} (430M params) as a frozen judge. DeBERTa's disentangled attention outperforms RoBERTa-large on MNLI by $\sim$1.2 points, which translates to roughly 90 fewer mislabels at our scale. Choosing a frozen NLI judge sidesteps both white-box requirements and LLM-as-judge bias, and lets the same judge be applied uniformly across all four models under test.

The main challenge is that NLI mislabels long correct answers as \lbl{neutral}: verbosity dilutes the entailment signal. We address this by (i) constraining the premise to gold supporting sentences only (never distractors), (ii) cleaning prompt prefixes and chat-template role markers from answers before NLI (otherwise the answer field contains the original prompt and ``\texttt{assistant\textbackslash n\textbackslash n}'' tokens), and (iii) layering the multi-signal override of Idea~3 on top.

\subsection{Idea 2: Score-Based Hallucination Type Classification}

Sequential heuristic chains fail on this task: whichever check fires first wins, even if a later check has stronger signal. Instead, each candidate type produces a score in $[0, 1]$ from independent signals, and we pick the argmax. Below a 0.5 confidence floor, we deliberately return \lbl{unsupported inference} rather than force a wrong type.

Each score is gated on a precondition that filters obvious false positives. \textit{Entity} score uses spaCy NER (PERSON/ORG/GPE/etc.) rather than capitalized-word matching, restricting to actual named entities. \textit{Attribute} score is gated on entity-correctness: number mismatches only count in answer sentences containing a confirmed-correct entity, so wrong-entity answers fall through to \lbl{entity error} instead. \textit{Multi-hop} score requires the answer to overlap at least 0.2 Jaccard with at least one evidence sentence---a model that invented an answer entirely is \lbl{unsupported inference}, not multi-hop. Before this redesign, \lbl{multi-hop reasoning error} had an estimated 80\% false-positive rate; after the gates and the score floor, the resulting type distribution carries real interpretive weight.

\subsection{Idea 3: Multi-Signal Supported Override and Cross-Record Detection}

The supported-override addresses NLI's verbose-correct failure mode. Three independent signals with uncorrelated failure modes compete: \textit{reference recall} (fraction of reference tokens in the answer; verbosity-tolerant), \textit{embedding cosine} (sentence-transformers/all-MiniLM-L6-v2 between answer and reference; captures paraphrase), and \textit{bidirectional NLI} (answer entails reference \textit{and} vice versa). Two of three must agree before flipping a record from hallucinated to supported. For \lbl{neutral}-NLI cases, reference recall above 0.5 alone suffices; for \lbl{contradiction}, we require the full 2-of-3 to avoid suppressing real errors. Thresholds (recall 0.50, cosine 0.60) were tuned on $\sim$1500 spot-checked records, biased toward false-negatives so we never artificially lower hallucination rates.

The \lbl{overconfident abstention failure} flag cannot be set per-record---it requires comparing a model's behavior on the same question across two prompt variants. We implement it as a cross-record second pass in \texttt{compute\_metrics.py}: group records by (model, question), and if \lbl{plain} is hallucinated and \lbl{abstain} is correctly abstained, flag the \lbl{plain} record. Isolating this in a separate pass keeps the per-record labeler stateless and the cross-record logic easy to audit.

\section{Experimental Setup (10 points)}

\subsection{Models}

Four open-weight instruction-tuned transformers, all run closed-book (question only, no context) at temperature 0.2, top-$p$ 0.9, max-new-tokens 256. Inference via Hugging Face \texttt{transformers} on Colab T4 (4-bit quantized) and A100 (FP16): \textbf{Phi-4-mini-instruct} (3.8B, Microsoft), \textbf{Mistral-7B-Instruct-v0.3} (7B, Mistral AI), \textbf{Qwen2.5-7B-Instruct} (7B, Alibaba), \textbf{Llama-3.1-8B-Instruct} (8B, Meta). The four span 3.8--8B and four provider families so observed differences reflect architecture/training rather than scale alone. Labeling additionally uses \texttt{microsoft/deberta-large-mnli} as the NLI judge, \texttt{sentence-transformers/all-MiniLM-L6-v2} for embeddings, and spaCy's \texttt{en\_core\_web\_sm} for NER.

\subsection{Dataset}

\textbf{HotpotQA (primary):} 500 questions sampled from the dev split (distractor config, seed 42), each with a short reference answer and human-annotated supporting facts (sentence pointers into gold passages). Models see only the question; supporting facts are reserved for the labeling pipeline. \textbf{TruthfulQA (secondary):} 120 questions stratified by category; lacks supporting facts, so the gold reference itself is used as the NLI premise. Test-set total: 500$\times$3 + 120$\times$3 = 1860 records per model, 7440 across all four. Prompt variants: \lbl{plain} (direct QA), \lbl{abstain} (``If unsure, say I don't know''), \lbl{reasoning} (step-by-step + abstention).

\subsection{Evaluation Metrics}

\textbf{Exact Match}: $\text{EM}(p, r) = \mathbb{1}[\text{normalize}(p) = \text{normalize}(r)]$ (lowercase, strip articles, drop punctuation).
\textbf{Token F1}: harmonic mean of token precision and recall: $F_1 = \frac{2 \cdot P \cdot R}{P + R}$.
\textbf{Hallucination rate}: fraction of records where \texttt{hallucination\_type} $\notin$ \{supported, abstained\}.
\textbf{Abstention rate}: fraction labeled \lbl{abstained}.
\textbf{Type distribution}: percent of hallucinated records per taxonomy bucket.
\textbf{Prompt sensitivity}: $\Delta = \text{Hall}\%_{\text{plain}} - \text{Hall}\%_{\text{reasoning}}$.
\textbf{Overconfident abstention failure}: cross-record fraction where \lbl{plain} hallucinates while \lbl{abstain} abstains. We additionally do a $\sim$25-example-per-type spot check (Section~\ref{sec:analysis}).

\section{Results (40 points for implementation and results)}
\label{sec:results}

\subsection{Hallucination, EM, F1 (HotpotQA)}

Table~\ref{tab:hall_rate} reports the four core metrics across all 12 (model $\times$ prompt) cells. EM is uniformly near zero because HotpotQA references are short strings while models generate full sentences; token-F1 is the more useful surface metric. Hallucination rate is the central measure.

\begin{table}[h!]
\centering
\small
\begin{tabular}{llrrrr}
\hline
Model & Prompt & EM & F1 & Hall\% & Abst\%\\
\hline
Phi-4-mini    & plain     & 0.008 & 0.066 & 87.6 & 0.0\\
Phi-4-mini    & abstain   & 0.008 & 0.017 & 3.0 & 92.2\\
Phi-4-mini    & reasoning & 0.000 & 0.036 & 81.8 & 6.6\\
Mistral-7B    & plain     & 0.002 & 0.048 & 92.4 & 0.2\\
Mistral-7B    & abstain   & 0.004 & 0.046 & 15.8 & 78.0\\
Mistral-7B    & reasoning & 0.000 & 0.021 & 75.4 & 5.4\\
Qwen2.5-7B    & plain     & 0.000 & 0.035 & 91.0 & 0.2\\
Qwen2.5-7B    & abstain   & 0.000 & 0.023 & 4.0 & 91.6\\
Qwen2.5-7B    & reasoning & 0.000 & 0.016 & 74.4 & 9.0\\
Llama-3.1-8B  & plain     & 0.004 & 0.069 & 40.2 & 0.4\\
Llama-3.1-8B  & abstain   & 0.008 & 0.024 & 6.0 & 88.4\\
Llama-3.1-8B  & reasoning & 0.000 & 0.032 & 58.0 & 0.8\\
\hline
\end{tabular}
\caption{HotpotQA results. EM and F1 against reference; Hall\% and Abst\% from our labeling pipeline. The \lbl{abstain} prompt drives near-total abstention in three of four models; only Mistral retains a meaningful answer rate.}
\label{tab:hall_rate}
\end{table}

\subsection{Hallucination Type Distribution}

\lbl{contradiction to evidence} dominates plain prompts (44--73\% across models): when models hallucinate, they typically do so confidently against the gold evidence, not merely uncertainly. Reasoning prompts redistribute errors: \lbl{attribute error} rises to 57\% for Mistral and 91.7\% for Qwen.

\begin{figure}[h!]
\centering
\includegraphics[width=\columnwidth]{plots/report_hallucination_rate_by_type.png}
\caption{Hallucination rate and type composition per model and prompt on HotpotQA. Bar height = total hallucination rate; segments show the share of each error type. Suffixes \texttt{P}/\texttt{R}/\texttt{A} indicate plain / reasoning / abstain prompts.}
\label{fig:hall_by_type}
\end{figure}

\subsection{Prompt Sensitivity and Cross-Dataset}

Table~\ref{tab:prompt_cross} combines prompt sensitivity on HotpotQA with the cross-dataset comparison. Qwen and Mistral drop $\sim$17 points under reasoning, Phi-4-mini drops 5.8, and Llama-3.1 \emph{rises} 17.8 points. The same prompt has opposite effects across model families. On TruthfulQA the four models converge near 83\% under reasoning, eliminating the model-specific HotpotQA pattern.

\begin{table}[h!]
\centering
\small
\begin{tabular}{lrrr|rr}
\hline
 & \multicolumn{3}{c|}{HotpotQA Hall\%} & \multicolumn{2}{c}{TruthfulQA Hall\%}\\
Model & plain & reason & $\Delta$ & plain & reason\\
\hline
Phi-4-mini    & 87.6 & 81.8 & $-$5.8           & 93.3 & 83.3\\
Mistral-7B    & 92.4 & 75.4 & $-$17.0          & 96.7 & 83.3\\
Qwen2.5-7B    & 91.0 & 74.4 & $-$\textbf{16.6} & 88.3 & 76.7\\
Llama-3.1-8B  & 40.2 & 58.0 & \textbf{+17.8}   & 92.5 & 83.3\\
\hline
\end{tabular}
\caption{Prompt sensitivity on HotpotQA (plain vs.\ reasoning, $\Delta$ in percentage points) and cross-dataset hallucination rates. TruthfulQA exhibits cross-model convergence near 83\% under reasoning that does not appear on HotpotQA.}
\label{tab:prompt_cross}
\end{table}

\subsection{Headline Findings}

(1) Aggregate rates are model-similar (87--92\%) on HotpotQA \lbl{plain}, with Llama-3.1 the outlier at 40\%. (2) Prompt sensitivity is sharply non-uniform: reasoning helps Mistral/Qwen, marginally helps Phi, hurts Llama. (3) Failure modes shift under reasoning: contradiction drops, attribute error rises---models become more careful about high-level facts but accumulate small numerical errors during elaboration. (4) Misconception questions defeat reasoning prompts uniformly across models on TruthfulQA.

\section{Analysis and Discussion (25 points)}
\label{sec:analysis}

We focus on \textbf{Mistral-7B}, the model with the clearest beneficial reasoning-prompt response (17-point drop) and the most balanced abstention behavior. Manual review of $\sim$25 hallucinations per type surfaces four recurring failure patterns: (1) \textit{confident substitution}---Mistral plain answers ``The Public Theater'' when asked the original name of Joseph Papp's organization (gold: New York Shakespeare Festival, the institution's earlier name); (2) \textit{over-elaboration with embedded fabrications}---Mistral reasoning correctly says ``Woolworth Building was built first'' but inserts ``completed in 1913'' (gold: 1910--1912); (3) \textit{confused chaining on bridge questions}---correct identification of an intermediate entity but wrong second-hop property; (4) \textit{hedged abstention on easy facts}---refusing questions whose evidence contains an unambiguous answer (e.g., D-Day for the Normandy operation).

\begin{figure}[h!]
\centering
\includegraphics[width=\columnwidth]{plots/report_type_shifts.png}
\caption{Mistral-7B hallucination type distribution shifts across prompts. Reasoning compresses contradiction errors into attribute errors; abstain shifts mass toward multi-hop and away from contradiction.}
\label{fig:type_shifts}
\end{figure}


\paragraph{H1: Reasoning prompts convert contradictions into attribute errors via verbosity.} Mistral's \lbl{contradiction to evidence} drops from 48.3\% to 11.4\% under reasoning while \lbl{attribute error} rises from 18.6\% to 57.0\%. Step-by-step reasoning improves high-level structure but lengthens the response, introducing extra numerical claims that miss the evidence (the Woolworth example above). Cap \texttt{max\_tokens} at 64 to test the verbosity mechanism.

\paragraph{H2: Bridge questions concentrate multi-hop errors.} HotpotQA's \texttt{bridge} questions require both identifying an intermediate entity \textit{and} retrieving a property of it---two failure points. Llama plain answers ``I'm not familiar with Pak Sao'' on a bridge question whose Wing-Chun hop is the bottleneck. The labeled records carry a \texttt{type} field; a short filter would separate bridge vs.\ comparison rates.

\paragraph{H3: Llama-3.1's anomalous prompt sensitivity is an instruction-following difference, not a knowledge difference.} Llama hallucinates 40\% on \lbl{plain} (vs.\ 87--92\% for the others) but \emph{rises} to 58\% on \lbl{reasoning}. Plain outputs frequently hedge (``I don't have information on...,'' caught by \texttt{ABSTAIN\_PHRASES}); reasoning outputs are multi-paragraph elaborations where the new hallucinations originate. Llama appears to read ``reason step-by-step'' as ``produce a longer answer'' rather than ``be more cautious.''

\paragraph{Pipeline limitations.} The spot-check surfaced two labeling weaknesses worth recording: (1) when gold evidence contains ambiguous referents (``the film features...'' detached from the article title), DeBERTa cannot attribute the sentence to its source entity, occasionally returning \lbl{neutral} on supported records; prepending each evidence sentence with its title is a natural fix. (2) \lbl{multi-hop reasoning error} retains some false-positive rate on verbose-but-correct answers; a stronger NER+coreference model or a GPT-4 tiebreaker on low-score cases would reduce it.

\begin{figure}[h!]
\centering
\includegraphics[width=\columnwidth]{plots/report_overconfident_abstention.png}
\caption{Overconfident abstention failures: questions where the plain prompt hallucinated but the abstain prompt correctly refused. Phi-4-mini and Qwen2.5 show the largest gaps between their plain and abstain behavior.}
\label{fig:overconfident}
\end{figure}


\paragraph{Takeaways.} Three points consolidate. First, prompt effects are mediated by \textit{response length}: reasoning prompts trade outright contradictions for embedded numerical errors. Second, prompt strategies that work on multi-hop factoid questions (HotpotQA) do not transfer to misconception-style questions (TruthfulQA), where all models converge near 83\%. Third, model-specific instruction-following style matters more than model size in determining how prompts shape hallucination behavior---hallucination-mitigation prompts must be model-conditioned rather than uniformly applied.

\section{Code and Demo (5 points)}

\paragraph{Repository.} \url{https://github.com/hpaliwal8/hallu-detection-type-classification} (TODO: confirm URL). Original code; third-party libraries are \texttt{transformers}, \texttt{datasets}, \texttt{sentence-transformers}, \texttt{spacy}, \texttt{seaborn}/\texttt{matplotlib}.

\paragraph{Key files.} Inference: \code{scripts/run_experiments.py}. Cleaning: \code{scripts/clean_outputs.py}. Labeling: \code{scripts/label_outputs.py} (calls \code{src/labeling/nli.py}, \code{src/labeling/heuristics.py} with score-based \code{classify_hallucination_type()} and the multi-signal \code{is_supported_by_signals()}, plus \code{src/labeling/embeddings.py}, \code{src/labeling/entities.py}, \code{src/labeling/evidence.py}). Metrics: \code{scripts/compute_metrics.py} with \code{flag_overconfident_abstention()}. Aggregation: \code{scripts/type_breakdown.py}. Plots: \code{scripts/plot_results.py}.

\paragraph{Reproducing results.} The end-to-end pipeline is six stages: prepare datasets, run inference per model, clean outputs, label, compute metrics, plot. Exact commands and per-stage flags are listed in the repository README; all commands take \texttt{--config config/default.yaml} and produce JSONL outputs under \texttt{data/outputs/}.

\paragraph{Models and data.} No models trained. Pre-trained checkpoints from Hugging Face: \texttt{microsoft/Phi-4-mini-instruct}, \texttt{mistralai/Mistral-7B-Instruct-v0.3}, \texttt{Qwen/Qwen2.5-7B-Instruct}, \texttt{meta-llama/Llama-3.1-8B-Instruct} (Meta license required), \texttt{microsoft/deberta-large-mnli}. Datasets: \texttt{hotpot\_qa} (distractor/validation) and \texttt{domenicrosati/TruthfulQA}. Prompts in \texttt{config/default.yaml}.

\paragraph{Software.} Python 3.12; \texttt{transformers}$\geq$4.40, \texttt{accelerate}$\geq$0.26, \texttt{datasets}$\geq$2.18, \texttt{sentence-transformers}$\geq$3.0, \texttt{spacy}$\geq$3.7 (\texttt{en\_core\_web\_sm}), \texttt{seaborn}$\geq$0.13, \texttt{matplotlib}$\geq$3.7, \texttt{pandas}$\geq$2.0, \texttt{bitsandbytes}$\geq$0.46.1 for T4 4-bit. T4 (16GB) sufficient with quantization; A100 preferred.

\paragraph{Demo.} (i) Show \code{compute_metrics.py} output matching Table~\ref{tab:hall_rate}; (ii) regenerate a plot via \code{plot_results.py}; (iii) walk a single record through cleaning $\rightarrow$ NLI $\rightarrow$ heuristic type assignment.

\section{Learning Outcomes}

\paragraph{Single classifiers need independent rescue paths.} Off-the-shelf DeBERTa-MNLI mislabels verbose-but-correct answers as \lbl{neutral}, and our multi-signal override (Idea 3) emerged only after manual false-positive review. The lesson generalizes: any classifier in a labeling pipeline benefits from an independent path that fires on its known weakness.

\paragraph{Score-based dispatch beats sequential heuristics; input cleaning beats both.} Our first \texttt{classify\_hallucination\_type} used \texttt{if/elif/else} and produced an estimated 80\% false-positive rate on \lbl{multi-hop reasoning error}; replacing the chain with score competition plus a 0.5 floor was a small refactor with disproportionate gains. But the largest single improvement of the whole project came from stripping prompt prefixes and chat-template role markers from answers via \texttt{clean\_outputs.py}---Mistral's supported count jumped from 119 to 416, larger than any subsequent algorithmic improvement. Audit pipeline inputs before tuning the model in the middle.

\paragraph{Aggregate metrics hide model differences; prompt sensitivity is an interaction effect.} Three of four models hallucinate at 87--92\% on HotpotQA \lbl{plain}---indistinguishable on a binary metric---but the type distribution reveals they fail differently. Similarly, ``add chain-of-thought instructions'' is not universally beneficial: the same reasoning prompt helps Qwen/Mistral by $\sim$17 points but hurts Llama-3.1 by 17.8 points. Comparative LLM evaluation should report fine-grained type distributions and per-model prompt sensitivity, not aggregate rates alone.

\section{Contributions}

\paragraph{Hitansh Paliwal.} TODO: list responsibilities.
\paragraph{Team Member 2 (TODO).} TODO: list responsibilities.

All members contributed to drafting this report. The entire team receives the same grade.

\end{document}
